\subsubsection{给定结构的连续精修与复核接口}
运输路线和载荷固定后，运输作业与充电时长为常数；中继位置固定后，转场时间为常数，服务时长决定线性悬停能耗。再固定各组件处于哪一充电分段，并保留该段能量上下界，资源接续就具有仿射形式。这些条件共同解释固定结构模型为什么可以采用线性规划。

分段边界是模型的一部分。只使用一段充电直线而不约束能量范围，可能算出原函数不支持的恢复时刻。两轮优化先压缩联合工期，再在给定工期容许量内减少不必要的服务能耗，所得解仍须返回完整物理计算和通信事件核验。Huangfu与Hall给出的对偶修正单纯形研究为所用线性规划技术提供方法背景\cite{huangfu2018}；本文的固定结构结论则来自明确的输入、约束及保存下界。

位置和任务组合选择、资源先后安排、具体开始时刻分属不同规模的决策。分层方法使每一层都有可列出的输入与输出，也使恢复计算能只调整受扰动的层次。数值保护间隔用于提高计算稳定性，执行保障范围则由本章实际压力情景和连续核验单独评价。
